\section{Darstellung der Ergebnisse}
Test 

Im folgenden Abschnitt wird die Umsetzung des entwickelten Prototyps beschrieben. 
Dabei werden zunächst die funktionalen und nicht-funktionalen Anforderungen dargestellt und danach die darauf basierende Softwarearchitektur sowie das Hardwarekonzept beschrieben. 

\subsubsection{Funktionale Anforderungen}

Das System muss über ein sprachbasiertes Ein- und Ausgabesystem verfügen, das gesprochene Anfragen erfassen und verarbeitete Antworten als Sprache ausgeben kann.

Eine Wakeword-Erkennung ist erforderlich, um das System gezielt aktivieren zu können, ohne dass eine dauerhafte Verarbeitung aller Audioeingaben stattfindet.

Darüber hinaus muss das System eine Erweiterbarkeit bieten, über die zusätzliche Funktionsmodule in das bestehende System integriert werden können.

Auf dem integrierten Display soll ein visueller Avatar dargestellt werden, der den aktuellen Systemzustand widerspiegelt und die Interaktion durch eine grafische Benutzeroberfläche ergänzt.

\subsubsection{Nicht-funktionale Anforderungen}

Die Bedienung des Systems muss intuitiv gestaltet sein, sodass keine technischen Vorkenntnisse für die Nutzung erforderlich sind.

Die verwendete Hardware muss kompakt dimensioniert sein, damit das Gesamtsystem als Schreibtischgerät eingesetzt werden kann.

Die Integration der einzelnen Hard- und Softwarekomponenten muss robust umgesetzt werden, um einen zuverlässigen Dauerbetrieb zu gewährleisten.

Das System soll echtzeitnahe Reaktionszeiten aufweisen, damit die sprachliche Interaktion als natürlich und flüssig wahrgenommen wird.

Die gesamte Interaktion, einschließlich der sprachlichen Ausgabe, muss verständlich und nachvollziehbar gestaltet sein.


\subsection{Softwarearchitektur}

Die Softwarearchitektur von Pablo basiert auf einer modularen, ereignisgesteuerten Python-Struktur mit Trennung zwischen Ein- und Ausgabe, Dialogverarbeitung sowie Hardwaresteuerung. 
Beim Start werden zunächst die Systemkonfiguration, die verfügbaren Plugin-Module und die visuelle Prozessdarstellung initialisiert. 
Anschließend wird der Wakeword- und Richtungsdetektor konfiguriert, der als Trigger-Komponente für die Interaktionsschleife fungiert.
Die aus der Wakeword-Analyse abgeleitete Schallrichtung wird unmittelbar auf die Servosteuerung übertragen, sodass sich der Bildschirm in Richtung der sprechenden Person ausrichtet.

Das System arbeitet mit einer Hauptschleife mit vier logisch getrennten Verarbeitungsphasen: 
(1) akustische Aktivierung, Richtungsschätzung der Stimme und Anpassung der Servoposition, 
(2) Spracherkennung mit Speech-to-Text (STT), 
(3) inhaltliche Verarbeitung über ein Large Language Model (LLM) mit optionaler Modulausführung und
(4) Sprachausgabe mit Text-to-Speech (TTS). 
Die einzelnen Phasen sind funktional entkoppelt und über definierte Funktionsaufrufe miteinander verbunden, wodurch die Austauschbarkeit einzelner Komponenten (z. B. STT- oder TTS-Backend) gewährleistet wird.

Ein zentrales Architekturelement ist die Plugin-Schnittstelle. Über den Loader werden externe Funktionsmodule dynamisch eingebunden und dem LLM-Kontext bereitgestellt. 
Durch den bereitgestellten Kontext kann das System zwischen generischer Antwortgenerierung und Modulausführung unterscheiden. 
Dies ermöglicht eine Erweiterung des Basissystem, ohne dass der Kernprozess verändert werden muss.

Parallelität wird gezielt für zeitkritische Teilaufgaben genutzt. 
Die Servosteuerung erfolgt asynchron in einem separaten Thread, sodass die Audioverarbeitung nicht blockiert wird. 
Zusätzlich wird die Zustandsumschaltung der Visualisierung während der Sprachausgabe zeitversetzt und threadbasiert umgesetzt, um eine konsistente Synchronisation zwischen auditivem und visuellem Feedback zu erreichen. 
Fehlerfälle werden über Ausnahmebehandlung auf Schleifenebene abgefangen. 

\subsubsection{(1) Aktivierung des Systems}

Die Aktivierung erfolgt über eine Wakeword-Erkennung mit Picovoice Porcupine. 
Beim Erkennen des Wakewords wird zeitgleich die Schallrichtung aus einem gepufferten Stereo-Audiofenster mittels GCC-PHAT-Korrelation der beiden Mikrofonkanäle bestimmt und zur Servo-Ausrichtung genutzt. 
Die Umsetzung erfolgt in einem integrierten Detector-Modul, das beide Aufgaben zusammenführt.


\subsubsection{(2) Spracherkennung (STT)}

Die Spracherkennung erfolgt mit dem Vosk-Framework. 
Der Audiostrom wird verarbeitet, bis das Vosk-Modell einen vollständigen Satz erkannt hat. 
Erst dann wird der Text an die Dialogverarbeitung übergeben.


\subsubsection{(3) LLM-Verarbeitung}

Die inhaltliche Verarbeitung nutzt eine API für den Zugriff auf ein Large-Language-Model (Llama 3.3 70B). 
Der erkannte Nutzereingabe-Text wird zusammen mit den vom Plugin-Loader geladenen Tool-Definitionen an das Modell übergeben. 
Das Modell entscheidet selbstständig (via \texttt{tool\_choice="auto"}), ob ein direkter Chat-Response ausreicht oder ob ein registriertes Plugin-Tool aufgerufen werden soll. 
Die Plugins kapseln die Logik für Wetter, Notizen, Kalender und Timer.


\subsubsection{(4) Sprachausgabe (TTS)}

Die Sprachausgabe wird lokal mit Piper generiert (deutsches Modell \texttt{de\_DE-thorsten-low.onnx}). 
Der Text wird in eine WAV-Datei gerendert und über einen verfügbaren System-Audio-Player (\texttt{paplay} oder \texttt{aplay}) abgespielt. 
Vor dem eigentlichen Audiomaterial wird eine führende Stille eingefügt, sodass die visuelle Feedback-Umschaltung mit dem tatsächlichen Audiobeginn zeitlich synchronisiert ist.


\subsection{Hardwarekonzept}
Die Auswahl der Hardwarekomponenten erfolgte unter Abwägung von Kosten, Kompaktheit, Integrationsfähigkeit und Community-Unterstützung.
Als zentrale Recheneinheit wurde ein Raspberry Pi gewählt, der aufgrund seiner kompakten Bauform, breiten Softwarekompatibilität und umfangreichen Community-Dokumentation eine geeignete Plattform für ein Embedded-Assistenzsystem darstellt.
Für die Audioerfassung wurde der ReSpeaker 2-Mics Pi HAT eingesetzt, der neben der Sprachaufnahme auch eine richtungsbasierte Schallquellenortung ermöglicht.
Die visuelle Ausgabe erfolgt über ein Waveshare 4-Zoll-IPS-Display, das direkt am Raspberry Pi betrieben wird.
Als Audioausgabe dient ein 5\,W / 8\,$\Omega$-Lautsprecher, der über einen separaten Audioverstärker angesteuert wird.
Die Stromversorgung des Verstärkers erfolgt über einen Arduino Uno, der zugleich als dedizierte Steuereinheit für die Energieversorgung der Audiokomponente fungiert.
Für die mechanische Ausrichtung des Kopfteils, das aus Display und Gehäuseoberteil besteht, wird ein Servomotor eingesetzt, der über einen Arduino Nano angesteuert und seriell vom Raspberry Pi kontrolliert wird.
Durch diese Aufteilung der Steuerungsaufgaben auf mehrere Mikrocontroller konnte die Echtzeitfähigkeit der mechanischen und audiovisuellen Komponenten sichergestellt werden, ohne die Hauptrecheneinheit zusätzlich zu belasten.

Das Gehäuse wurde als 3D-gedruckter Baumstamm gestaltet und nimmt sämtliche Hardwarekomponenten auf.
Diese Formgebung unterstreicht den Companion-Charakter von Pablo und schafft eine visuelle Abgrenzung gegenüber rein technischen Geräten.